\section{Related Work}
\label{sec:related}

% need to validate citation
% need to add more literatures
The development and evaluation of machine learning models for cybersecurity, particularly for intrusion detection and anomaly detection, critically rely on realistic and comprehensive datasets. Historically, benchmarks like DARPA 1998 and KDD Cup 1999~\cite{tavallaee2009detailed} served this purpose but became outdated due to evolving attack methods and simplistic traffic patterns. More recent datasets such as UNSW-NB15~\cite{moustafa2015unsw} and CICIDS2017~\cite{sharafaldin2018toward} introduced modern attack vectors and richer network traffic characteristics. However, these datasets remain static, lacking the flexibility to control factors critical for robustness evaluation, such as the prevalence of specific attack types, temporal dynamics, or data scarcity scenarios. This static nature limits their utility for systematic sensitivity analysis and benchmarking under diverse operational conditions and evolving threat landscapes.

To address the limitations of static real-world datasets, particularly challenges like data scarcity, privacy concerns, difficulty in capturing rare or emerging threats, and the lack of control over data distributions, synthetic data generation has emerged as a promising avenue. Existing approaches for generating synthetic cybersecurity data span various techniques. Early methods often relied on statistical models or rule-based systems to mimic traffic patterns. More recently, deep generative models, especially Generative Adversarial Networks (GANs), have shown significant potential for generating tabular network traffic data with high fidelity and utility. Ammara et al.~\cite{ammara2025syntheticnetworktrafficdata} conducted a comparative analysis of statistical, conventional AI, and GAN-based methods (like CTGAN and CopulaGAN), highlighting GANs' superiority in capturing complex data distributions and handling class imbalance in datasets like NSL-KDD and CICIDS-2017. Similarly, Zouhri and Idri~\cite{10765501} specifically applied CTGAN to address multi-class imbalance in CIC-IDS2018, demonstrating improved detection performance compared to traditional oversampling techniques. Rahman et al.~\cite{rahman2025leveraging} proposed integrating GANs with classifiers like Random Forest to generate synthetic attack samples, enhancing dataset diversity and improving the detection of minority attack classes, thus improving classifier generalization and robustness.

Beyond tabular traffic data, Large Language Models (LLMs) are increasingly explored for generating diverse forms of cybersecurity data, including security logs, attack scenarios, and domain-specific text. Razmyslovich et al.~\cite{razmyslovich2025eltex} introduced ELTEX, a framework leveraging LLMs for domain-driven synthetic text generation, demonstrating its effectiveness for cybersecurity text classification by integrating domain indicators. However, their application to structured \textit{tabular} data is still nascent and faces considerable hurdles. The synthetic data generated by these methods is critical for tackling specific challenges in cybersecurity ML, such as class imbalance and data scarcity. Class imbalance, where attack instances are significantly outnumbered by benign traffic, is a pervasive issue leading to biased model performance~\cite{5356528}. Research by Lian et al.~\cite{lian2021robustness} demonstrates the importance of evaluating classifier robustness under varying imbalance conditions using a design-of-experiments approach. Sethi et al.~\cite{deepcyber2024imbalance} explored deep learning techniques using imbalance-aware training methods like oversampling and cost-sensitive learning in cybersecurity contexts. The survey by Lu et al.~\cite{Chen2024ASO} provides a broader overview of imbalanced learning techniques. While synthetic data generation techniques (like GANs and LLMs) are used to create data that addresses these issues (e.g., oversampling minority classes), there is a distinct need for frameworks that allow researchers to systematically control and vary the degree and nature of imbalance, data scarcity, and rare event characteristics within generated datasets. This fine-grained control is essential for conducting rigorous, reproducible experiments to evaluate model sensitivity, robustness, and generalization capabilities across a spectrum of challenging real-world scenarios.

% Existing synthetic data methods primarily focus on generating realistic data distributions or augmenting existing minority classes to improve training data quantity or balance. However, they often lack the explicit, configurable control required to generate datasets with precisely specified characteristics regarding attack prevalence, temporal patterns, semantic attributes, and structured event narratives, tailored for systematic evaluation. This limitation hinders researchers' ability to conduct controlled experiments to understand why a model performs well or poorly under specific, challenging data conditions, especially concerning rare attacks and varying levels of data scarcity or imbalance. The field currently lacks a comprehensive, configurable framework designed specifically for generating synthetic cybersecurity datasets tailored for rigorous, systematic evaluation of model sensitivity, robustness, and interpretability under a wide range of user-defined conditions.

% limitation and gap
While deep generative models, particularly Generative Adversarial Networks (GANs), have advanced synthetic cybersecurity data generation beyond traditional statistical methods~\cite{ammara2025syntheticnetworktrafficdata}, they are not without significant limitations. GANs, including variants like CTGAN and WGAN, often suffer from training instability and mode collapse, which can lead to the generation of monotonous or low-diversity attack samples, failing to capture the full spectrum of novel and complex threats~\cite{bourou2021review}. General-purpose LLMs often fail to maintain the strict logical and statistical relationships between columns in cybersecurity datasets (e.g., consistency between port numbers and services), a challenge highlighted in recent surveys on synthetic tabular data generation~\cite{gao2025comprehensivesurveyimbalanceddata}. These models are primarily designed for data augmentation or distribution replication rather than for creating datasets with precisely controlled characteristics. Crucially, the predominant focus of existing research—whether GAN or LLM-based—has been on achieving high-fidelity data replication for a single data distribution. This leaves a critical gap: the lack of a configurable framework designed for the systematic benchmarking of model robustness. The field lacks tools that allow researchers to meticulously control and vary dataset parameters such as attack prevalence, class imbalance ratios, temporal event dynamics, and the introduction of rare attack signatures. This fine-grained control is essential for conducting rigorous sensitivity analysis and understanding model behavior under the challenging and dynamic conditions of real-world deployments~\cite{lian2021robustness}.

Addressing this critical gap, our proposed CyberDataGen framework offers a novel approach to synthetic cybersecurity dataset generation. Unlike static benchmarks or general-purpose data augmentation techniques, CyberDataGen is engineered to provide fine-grained control over dataset characteristics, including variable attack ratios, prevalence of rare events, temporal dynamics, and rich semantic labeling integrated with structured event narratives. This configurability enables researchers to systematically generate a multitude of distinct datasets from a single framework, allowing for robust sensitivity analysis and benchmarking of ML model performance and behavior under user-defined data scarcity, class imbalance, and rare event scenarios. Furthermore, by generating data with structured semantic annotations and integrating potential LLM-based components for event narratives, CyberDataGen facilitates enhanced explainability studies using techniques like SHAP and LIME, providing insights into model decision-making processes for specific synthetic events. This combination of configurable realism, systematic evaluation capabilities, and support for explainability sets CyberDataGen apart and establishes a new standard for benchmarking cybersecurity ML systems.